% !TeX root = ../../Book2.tex
% 纲要标题：### 11.4 中心化子
% 纲要位置：计划纲要.md，第 1426 行。

\section{中心化子}
\label{b2:sec:1104}

一个作用确定以后，可以再问哪些线性算子与它交换。这些算子组成中心化子；再取一次中心化子，则至少能找回原来的作用。但是否恰好找回原代数，是一个需要证明的问题。本节先给出完全可算的例子，再说明它与张量表示的联系。

% 纲要要点（供目录核对）：
%
% * 模自同态与中心化子，双中心化子的包含与幂等性
% * 正则左右作用及有限维张量因子的中心化子
% * 三角反例与 Schur–Weyl 对偶的动机及证明归属
%
% ---
%

\subsection{交换作用与中心化子}
\label{b2:subsec:1104:01}

设 \(k\) 为域、\(V\) 为 \(k\) 向量空间，\(\mathcal A\subseteq\operatorname{End}_k(V)\) 为含恒等算子的子代数。沿用\subsecref{b2:subsec:0804:03}的中心化子定义，记
\[
\mathcal A'=\bigl\{\,T\in\operatorname{End}_k(V):Ta=aT\text{ 对所有 }a\in\mathcal A\,\bigr\}.
\]
再记 \(\mathcal A''=(\mathcal A')'\)，称为\term[shuang-zhongxinhua-zi]{双中心化子}[double centralizer]。这些撇号在本节表示中心化子，不表示导数或对偶。
\symidx{\(\mathcal A'\)、\(\mathcal A''\)：给定作用的中心化子及双中心化子}

\begin{proposition}\label{b2:prop:commutant-endomorphism}
设 \(k\) 为域、\(A\) 为结合含幺 \(k\) 代数、\(V\) 为 \(k\) 向量空间，\(\rho:A\to\operatorname{End}_k(V)\) 为含幺表示。对 \(\mathcal C\subseteq\operatorname{End}_k(V)\)，用 \(\mathcal C'\) 表示 \(\operatorname{End}_k(V)\) 中与 \(\mathcal C\) 全部元素交换的算子所成的中心化子，则
\[
\rho(A)'=\operatorname{End}_A(V),\qquad \rho(A)\subseteq\rho(A)''.
\]
若 \(\mathcal A\subseteq\mathcal B\subseteq\operatorname{End}_k(V)\)，则
\[
\mathcal B'\subseteq\mathcal A',\qquad
\mathcal A''\subseteq\mathcal B''.
\]
对任意算子族 \(\mathcal C\subseteq\operatorname{End}_k(V)\)，还有
\[
\mathcal C\subseteq\mathcal C'',\qquad
\mathcal C'''=\mathcal C',\qquad
\mathcal C''''=\mathcal C''.
\]
\end{proposition}
\begin{proof}
\cref{b2:prop:algebra-modules-representations}中模同态的相容公式在定义域、陪域均为 \(V\) 时，就是 \(T\rho(a)=\rho(a)T\)。在向量 \(v\) 上，这个等式写成 \(T(av)=a\cdot T(v)\)，故第一个等式成立。任意 \(c\in\mathcal C\) 都与 \(\mathcal C'\) 的全部元素交换，所以 \(\mathcal C\subseteq\mathcal C''\)，这也给出 \(\rho(A)\subseteq\rho(A)''\)。

若一个算子与 \(\mathcal B\) 的全部元素交换，它当然与子集 \(\mathcal A\) 的全部元素交换，故 \(\mathcal B'\subseteq\mathcal A'\)。再取一次中心化子并反转包含，得到 \(\mathcal A''\subseteq\mathcal B''\)。由 \(\mathcal C\subseteq\mathcal C''\) 得 \(\mathcal C'''\subseteq\mathcal C'\)；把扩大性质用于 \(\mathcal C'\)，又得 \(\mathcal C'\subseteq\mathcal C'''\)。于是 \(\mathcal C'''=\mathcal C'\)，两侧再取中心化子即得 \(\mathcal C''''=\mathcal C''\)。
\end{proof}

与更多算子交换会施加更多限制，所以中心化子反转包含；连续取两次后又保持包含。上述等式还说明，双中心化子扩大原算子族，并且再取双中心化子不会继续增大。这是它作为一种闭包运算的意义。

\ProseParagraphBreak
若两个代数在同一空间上的作用彼此交换，其中一个便落在另一个的中心化子中。只有证明反向包含，才能称它们互为中心化子。维数有限时可以把交换条件写成线性方程组；这通常比直接猜测所有相容映射更可靠。

\subsection{双中心化子的基本例子}
\label{b2:subsec:1104:02}

\begin{theorem}[正则作用的双中心化子]\label{b2:thm:regular-double-centralizer}
设 \(k\) 为域、\(A\) 为结合含幺 \(k\) 代数。在 \(\operatorname{End}_k(A)\) 中，令
\[
L(A)=\{L_a:x\mapsto ax\mid a\in A\},\qquad
R(A)=\{R_a:x\mapsto xa\mid a\in A\}.
\]
映射 \(a\mapsto L_a\) 与 \(a^{\mathrm{op}}\mapsto R_a\) 分别给出 \(k\) 代数同构 \(A\cong L(A)\) 与 \(A^{\mathrm{op}}\cong R(A)\)。两个作用代数互为中心化子，因此 \(L(A)''=L(A)\)，不要求 \(A\) 有限维。
\end{theorem}
\begin{proof}
复合满足 \(L_aL_b=L_{ab}\)、\(R_aR_b=R_{ba}\)，所以右乘的复合顺序对应反代数乘法。两种映射都线性、保单位元，并且算子在 \(1\) 处的值恢复 \(a\)，故它们单射；由像的定义，它们满射到相应作用代数。

\(T\) 与全部左乘交换，当且仅当它左 \(A\) 线性。正则左模由 \(1\) 生成；\cref{b2:prop:regular-bimodule-endomorphism}给出这样的 \(T\) 由 \(b=T(1)\) 唯一确定，且 \(T(x)=xb\)，故 \(L(A)'=R(A)\)。若 \(T\) 与全部右乘交换，则对任意 \(a\)，
\[
T(a)=T\bigl(R_a(1)\bigr)=R_a\bigl(T(1)\bigr)=T(1)a.
\]
整个算子仍由 \(T(1)\) 决定，这次它是左乘 \(T(1)\)。反向左乘与右乘由结合律交换，得到 \(R(A)'=L(A)\)。
\end{proof}

另一个重要例子来自两个向量空间的张量。令 \(U,W\) 为非零有限维空间。代数 \(\operatorname{End}_k(U)\) 作用于第一因子，\(\operatorname{End}_k(W)\) 作用于第二因子，两种作用交换。选定 \(U\) 的一组 \(m\) 元基后，\(U\otimes_kW\) 是 \(W\) 的 \(m\) 个副本的直和；第一因子的矩阵作用调配这些副本。与所有这些调配交换的算子，必须在每个副本上作同一个变换。

\begin{proposition}\label{b2:prop:tensor-factor-centralizers}
设 \(k\) 为域，\(U,W\) 为非零有限维 \(k\) 向量空间。在 \(\operatorname{End}_k(U\otimes_kW)\) 中，
\[
\operatorname{End}_k(U)\otimes k\operatorname{id}_W,
\qquad k\operatorname{id}_U\otimes\operatorname{End}_k(W)
\]
互为中心化子。此外，映射
\[
\begin{aligned}
&\operatorname{End}_k(U)\otimes_k\operatorname{End}_k(W)
\longrightarrow\operatorname{End}_k(U\otimes_kW),\\
&F\otimes G\longmapsto
\bigl(u\otimes w\mapsto F(u)\otimes G(w)\bigr)
\end{aligned}
\]
是 \(k\) 代数同构。
\end{proposition}
\begin{proof}
选 \(U\) 的基 \(u_1,\ldots,u_m\)。任意线性算子 \(T\) 可唯一写成
\[
T(u_j\otimes w)=\sum_i u_i\otimes T_{ij}(w),
\qquad T_{ij}\in\operatorname{End}_k(W).
\]
\(E_{aa}\otimes\operatorname{id}_W\) 是到第 \(a\) 个副本 \(ku_a\otimes W\) 的投影。若 \(T\) 与这些投影交换，它就保留各个副本，故 \(T_{ij}=0\) 当 \(i\ne j\)。而 \(E_{ab}\otimes\operatorname{id}_W\) 把第 \(b\) 个副本移到第 \(a\) 个副本，在其中保留 \(w\) 不变。与这个算子交换，便要求两个副本上的变换相同：在 \(u_b\otimes w\) 上比较，得到 \(T_{aa}(w)=T_{bb}(w)\)。因此 \(T=\operatorname{id}_U\otimes S\)，其中 \(S\) 为公共对角块。反向这样的算子确实交换，得到第一个中心化子。交换两因子的角色，得到另一个。

再选 \(W\) 的基 \(w_1,\ldots,w_\ell\)。张量 \(E_{ij}\otimes E_{ab}\) 将基向量 \(u_j\otimes w_b\) 送到 \(u_i\otimes w_a\)，将其余基向量送到零，所以这些张量恰对应全部矩阵单位。它们在两边均为基，故作用映射是线性同构；逐因子复合保证它保持乘法与单位。
\end{proof}

例如 \(k^n\oplus k^n\cong k^n\otimes_k k^2\)（\(n\ge1\)）时，\(\operatorname{diag}(B,B)\)、\(B\in M_n(k)\)，作用于第一因子。其中心化子为 \(\operatorname{id}_{k^n}\otimes\operatorname{End}_k(k^2)\)，允许在两个相同的 \(k^n\) 副本之间任意混合，而在各副本内部保留原来的矩阵作用。\cref{b2:ex:11-repeated-matrix-blocks}将这个结论写成具体的分块矩阵。

\ProseParagraphBreak
这里自同态代数的张量同构不能无条件推广到无限维空间。取 \(U,W\) 分别以 \(u_1,u_2,\ldots\) 与 \(w_1,w_2,\ldots\) 为基，定义
\[
Q(u_i\otimes w_j)=
\begin{cases}
u_i\otimes w_i,&i=j,\\
0,&i\ne j.
\end{cases}
\]
这在张量积的一组基上规定了像，故给出线性算子。若 \(Q=\sum_{r=1}^s A_r\otimes B_r\)，记 \(\alpha_i:U\to k\) 为第 \(i\) 个坐标泛函，对输入 \(u_i\otimes w\) 的输出取第 \(i\) 个 \(U\) 坐标，就得到
\[
P_i(w)=\sum_{r=1}^s\alpha_i(A_ru_i)B_r(w),
\]
其中 \(P_i\in\operatorname{End}_k(W)\) 保留 \(w_i\) 坐标，并把其余坐标归零。因此所有 \(P_i\) 都在有限维空间 \(\operatorname{span}_k\{B_1,\ldots,B_s\}\) 中。但这些投影线性无关：任意有限线性关系分别作用于对应的 \(w_i\)，就迫使每个系数为零。矛盾。所以这个 \(Q\) 不在张量作用映射的像中。

\ProseParagraphBreak
双中心化子等式不能推广到每个子代数。取 \(\mathcal A=T_2(k)\subseteq M_2(k)\)。与 \(E_{11}\) 交换迫使矩阵为对角矩阵，再与 \(E_{12}\) 交换迫使两个对角元相等。因此 \(\mathcal A'=kI_2\)，而 \(\mathcal A''=M_2(k)\supsetneq\mathcal A\)。原代数保留非零真子空间 \(ke_1\)，故作用不是不可约的。这就给出了从中心化子仅含标量反推不可约性的反例。

\subsection{Schur–Weyl 对偶的环论背景}
\label{b2:subsec:1104:03}

\chapref{b2:ch:08}已在 \(V^{\otimes d}\) 上构造一般线性群的对角作用与 \(S_d\) 的置换作用，并在\cref{b2:prop:commuting-tensor-actions}中证明它们交换。因而两种作用的线性生成代数互相包含于对方的中心化子。对于有限维复空间，\cref{b2:claim:schur-weyl-preview}给出它们互为中心化子的等式；证明见第五卷第7.3节。

\ProseParagraphBreak
这里的情形与\cref{b2:prop:tensor-factor-centralizers}不同。后者允许独立选取 \(F\in\operatorname{End}_k(U)\) 与 \(G\in\operatorname{End}_k(W)\)，作用为 \(u\otimes w\mapsto F(u)\otimes G(w)\)；Schur--Weyl 情形则让同一个 \(g\) 同时作用在所有因子上，得到 \(v_1\otimes\cdots\otimes v_d\mapsto gv_1\otimes\cdots\otimes gv_d\)，另一个作用改变因子的排列。独立因子上的中心化子计算没有计算出同一个 \(g\) 的对角作用的全部中心化子，因而不能代替 Schur--Weyl 对偶中反向包含的证明。

\ProseParagraphBreak
若只考察交换两个因子的算子 \(\tau\)，在特征不为二的域上，第八章的两个投影 \((\operatorname{id}\pm\tau)/2\) 把张量积分成对称与交替两部分。与 \(\tau\) 交换的算子恰好分别保持这两部分，因而可以在每一部分内部独立作用。\cref{b2:ex:11-tensor-flip-centralizer}在 \(V=k^2\) 上计算这一中心化子及其双中心化子。

\ProseParagraphBreak
\chapref{b2:ch:14}将研究不可约作用的中心化子及稠密性；有限维且底域代数闭时，不可约性会迫使作用代数等于整个自同态代数。上三角例子则表明，只知中心化子由标量组成还不足以作这个推断。双中心化子问题的答案取决于整个作用，而不只取决于原代数的抽象乘法表。

\ProseParagraphBreak
若另具备Hilbert空间和有界算子的基础，可参阅\appref{b2:app:D}的\cref{b2:thm:operator-bicommutant}。该结果在含幺且对伴随封闭的条件下，把双交换子与强、弱算子闭包联系起来；这里的纯代数中心化子计算并不自动包含那些拓扑条件。
